\begin{textsample}{POS Dim 2 – to – Score 23.00 – to/to\_em\_41.txt}  \label{ex:f2_pos_135}
3.3.1 \textbf{Carga} \textbf{Horária} Como já foi referenciado, o Novo Ensino Médio trouxe mudanças na \textbf{carga} \textbf{horária} dos \textbf{cursos}, promovendo um aumento progressivo da \textbf{carga} \textbf{horária} no Ensino Médio de 800 para 1.000 horas por ano letivo ( 3.000 horas em todo o Ensino Médio ).
Das 3.000 horas totais do Ensino Médio, fica estabelecido que 1.800 devam ser dedicadas ao cumprimento da BNCC, enquanto às outras 1.200 horas serão dedicadas aos \textbf{itinerários} \textbf{formativos}.
A \textbf{carga} \textbf{horária} mínima dos \textbf{cursos} \textbf{técnicos} da Educação Profissional Integrado ao Ensino Médio com \textbf{oferta} na forma articulada com o Ensino Médio — integrada ou concomitante — é indicada no \textbf{Catálogo} Nacional de \textbf{Cursos} \textbf{Técnicos}.
As \textbf{habilitações} \textbf{profissionais} técnicas podem ter \textbf{carga} \textbf{horária} mínima variando entre 800, 1.000 e 1.200 horas, e os \textbf{cursos} de \textbf{qualificação} \textbf{técnicas} — FIC — \textbf{cargas} \textbf{horárias} entre 160, 200, 240 e 400 horas. a ) Possibilidades de arquiteturas curriculares do Itinerário da Formação \textbf{Técnica} e \textbf{Profissional} com forma de \textbf{oferta} integrado ao Ensino Médio no \textbf{regime} parcial, respeitando a \textbf{carga} \textbf{horária} de 3.000 horas.
Atenção : Os exemplos abaixo serão organizados considerando o tempo de duração efetivo de aula de 50 minutos.
Para isso, foram realizados os cálculos para conversão das \textbf{cargas} \textbf{horárias} de hora ( h ) para hora/aulas ( h/a ) conforme cronograma abaixo : | Hora | Hora/aula ( 50 min ) | |---:|---:| | 1.800 horas | 2.160 horas/aula | | 1.200 horas | 1.440 horas/aula | | 1.000 horas | 1.200 horas/aula | | 800 horas | 960 horas/aula | | 3.000 horas | 3.600 horas/aula | | 4.500 horas | 5.400 horas/aula | 1 ) \textbf{Curso} de \textbf{habilitação} \textbf{técnica} com \textbf{carga} \textbf{horária} mínima obrigatória de 800 horas ( 960 horas/aulas ) integrado ao Ensino Médio.
Essa proposta de arquitetura apresenta a seguinte flexibilização : Formação Geral Básica de 2.160 horas/aulas ; formação \textbf{técnica} e \textbf{profissional} de 960 horas/aulas ; projeto de vida de 240 horas/aulas e \textbf{eletivas}, de 240 horas/aulas. | Série | Formação Geral Básica | Formação \textbf{Técnica} e Profissional | Projeto de Vida | \textbf{Eletivas} | \textbf{Carga} \textbf{Horária} ( h/a ) | |---|---:|---:|---:|---:|---:| | 1ª | 720 | 320 | 80 | 80 | 1.200 | | 2ª | 720 | 320 | 80 | 80 | 1.200 | | 3ª | 720 | 320 | 80 | 80 | 1.200 | | Total | 2.160 | 960 | 240 | 240 | 3.600 h/a | 2 ) \textbf{Curso} de \textbf{Habilitação} \textbf{técnica} de 1.000 horas ( 1.200 horas/aulas ) integrado ao Ensino Médio. | Série | Formação Geral Básica | Formação \textbf{Técnica} e Profissional | Projeto de Vida | \textbf{Eletivas} | \textbf{Carga} \textbf{Horária} ( h/a ) | |---|---:|---:|---:|---:|---:| | 1ª | 720 | 400 | 40 | 40 | 1.200 | | 2ª | 720 | 400 | 40 | 40 | 1.200 | | 3ª | 720 | 400 | 40 | 40 | 1.200 | | Total | 2.160 | 1.200 | 120 | 120 | 3.600 h/a | Essa possibilidade de arquitetura apresenta a Formação Geral Básica com 2.160 horas/aulas, formação \textbf{técnica} e \textbf{profissional} de 1.200 horas/aulas, projeto de vida com 120 horas/aulas e \textbf{eletivas} com 120 horas/aulas. 3 ) \textbf{Curso} de \textbf{Habilitação} \textbf{técnica} de 1.200 horas ( 1.440 horas/aulas ) integrado ao Ensino Médio Essa proposta traz uma arquitetura curricular com 2.160 horas/aulas dedicadas a Formação Geral Básica, 1.440 horas/aulas de formação \textbf{técnica} e \textbf{profissional}.
Nesse cenário, o projeto de vida será trabalhado de forma interdisciplinar, trabalhando competências que estimulem a autonomia do estudante, favorecendo a busca de conhecimento e as habilidades necessárias para tomar suas decisões. 3.4 \textbf{Curso} de \textbf{Habilitação} \textbf{técnica} de 1.200 horas ( 1.440 horas/aulas ) integrado ao Ensino Médio | Série | Formação Geral Básica | Formação \textbf{Técnica} e Profissional | \textbf{Carga} \textbf{Horária} ( h/a ) | |---|---:|---:|---:| | 1ª | 720 | 480 | 1.200 | | 2ª | 720 | 480 | 1.200 | | 3ª | 720 | 480 | 1.200 | | Total | 2.160 | 1.440 | 3.600 h/a | | Série | Formação Geral Básica | Formação \textbf{Técnica} e Profissional | Projeto de Vida | \textbf{Carga} \textbf{Horária} ( h/a ) | |---|---:|---:|---:|---:| | 1ª | 720 | 480 | 40 | 1.240 | | 2ª | 720 | 480 | 40 | 1.240 | | 3ª | 720 | 480 | 40 | 1.240 | | Total | 2.160 | 1.440 | 120 | 3.720 h/a | Essa segunda proposta de arquitetura curricular contempla 2.160 horas/aulas de Formação Geral Básica, 1.440 horas/aulas de formação \textbf{técnica} e \textbf{profissional} e 120 horas/aulas de projeto vida.
Essa sugestão considera trabalhar uma \textbf{carga} \textbf{horária} que exceda as 1.000 horas anuais, ficando estabelecido o cumprimento de 20 \% da \textbf{carga} \textbf{horária} total em atividades não presenciais conforme assegurado pela \textbf{legislação} \textbf{vigente}.
Possibilidades de arquiteturas curriculares do Itinerário da Formação \textbf{Técnica} e \textbf{Profissional} integrado ao Ensino Médio, como \textbf{oferta} no \textbf{regime} integral com cumprimento de \textbf{carga} \textbf{horária} de 4.500 horas ( 5.400 horas/aulas ) : 1 ) \textbf{Curso} de \textbf{habilitação} \textbf{técnica} de 800 horas ( 960 horas/aulas ). | Série | Formação Geral Básica | Formação \textbf{Técnica} e Profissional | Projeto de Vida | \textbf{Eletivas} | Trilha de Aprofundamento das áreas | Unidades Curriculares Integradoras | \textbf{Carga} \textbf{Horária} ( h/a ) | |---|---:|---:|---:|---:|---:|---:|---:| | 1ª | 720 | 320 | 80 | 120 | 120 | 440 | 1.800 | | 2ª | 720 | 320 | 80 | 120 | 120 | 440 | 1.800 | | 3ª | 720 | 320 | 80 | 120 | 120 | 440 | 1.800 | | Total | 2.160 | 960 | 240 | 360 | 360 | 1.320 | 5.400 h/a | Nessa arquitetura, há a proposição de integração de uma Trilha de Aprofundamento das áreas do conhecimento de 360 horas/aulas com o \textbf{itinerário} da formação \textbf{técnica} e \textbf{profissional}. 2 ) \textbf{Curso} de \textbf{habilitação} \textbf{técnica} de 1.000 horas ( 1.200 horas/aulas ). | Série | Formação Geral Básica | Formação \textbf{Técnica} e Profissional | Projeto de Vida | \textbf{Eletivas} | Unidades Curriculares Integradoras | \textbf{Carga} \textbf{Horária} ( h/a ) | |---|---:|---:|---:|---:|---:|---:| | 1ª | 720 | 400 | 80 | 120 | 480 | 1.800 | | 2ª | 720 | 400 | 80 | 120 | 480 | 1.800 | | 3ª | 720 | 400 | 80 | 120 | 480 | 1.800 | | Total | 2.160 | 1.200 | 240 | 360 | 1.440 | 5.400 h/a | 3 ) \textbf{Curso} de \textbf{habilitação} \textbf{técnica} de 1.200 horas ( 1.440 horas/aulas ). | Série | Formação Geral Básica | Formação \textbf{Técnica} e Profissional | Projeto de Vida | \textbf{Eletivas} | Unidades Curriculares Integradoras | \textbf{Carga} \textbf{Horária} ( h/a ) | |---|---:|---:|---:|---:|---:|---:| | 1ª | 720 | 480 | 80 | 120 | 400 | 1.800 | | 2ª | 720 | 480 | 80 | 120 | 400 | 1.800 | Com base nas inovações trazidas pela \textbf{Reforma} do Ensino Médio, no território do Tocantins é possível planejar uma organização curricular que promova integração de Trilhas de Aprofundamento das áreas do conhecimento com \textbf{itinerário} \textbf{formativo} da formação \textbf{técnica} e \textbf{profissional} ( \textbf{qualificação} \textbf{profissional} ou \textbf{habilitação} \textbf{técnica} \textbf{profissional} ) em sua parte \textbf{flexível}.
Esse cenário deve levar em consideração as áreas de interesse em que os estudantes querem aprofundar os conhecimentos, uma vez que essas unidades curriculares têm o objetivo de contribuir para o desenvolvimento de competências e habilidades fundamentais para a \textbf{trajetória} \textbf{profissional}, pessoal e social dos estudantes. 4 ) \textbf{PERFIL} DOCENTE Para o território do Tocantins, nos \textbf{cursos} pertencentes a modalidade da Educação Profissional \textbf{Técnica} de Nível Médio, os docentes devem ter formação de nível superior, em \textbf{curso} de licenciatura, \textbf{cursos} de bacharelado ou \textbf{cursos} superiores tecnólogo realizados em universidades ou demais instituições superiores de educação.
Esses critérios devem ser observados no instante da modulação desses sujeitos, relacionando a área atuação com a formação acadêmica.
A práxis de professores e/ou tutores que atuarão com o \textbf{itinerário} de formação \textbf{técnica} e \textbf{profissional} deve estar associada a uma postura mais dinâmica, capaz de proporcionar aos estudantes modelos de aprendizagem que os habilite para uma vida pessoal, social e \textbf{profissional} \textbf{alinhada} aos seus projetos de vida e as demandas da sociedade.
Para isso, é importante que o docente contemple no seu planejamento a articulação entre a teoria e a aplicação prática, associando -as ao contexto \textbf{profissional} optado, como forma de ampliar a visão de mundo e garantir a formação integral do estudante.
A \textbf{Reforma} do Ensino Médio também trouxe alterações a educação \textbf{profissional}, permitindo que \textbf{profissionais} com \textbf{cursos} de complementação pedagógica ou considerados com \textbf{notório} saber possam ministrar aulas, sem a necessidade de uma \textbf{certificação} de licenciatura.
Os procedimentos a serem adotados para o reconhecimento de \textbf{profissionais} com \textbf{notório} saber será \textbf{normatizado} pelo Conselho Estadual de Educação do Tocantins, com fundamento ao disposto no inciso V do caput do artigo 36 da LDB com redação alterada pela Lei nº 13.415/2017.

% matched lemmas: alinhar, carga, catálogo, certificação, curso, eletivo, flexível, formativo, habilitação, horário, itinerário, legislação, normatizar, notório, oferta, perfil, profissional, qualificação, reforma, regime, trajetória, técnico, vigente
\end{textsample}
